% ============================================================
%  Section 5: Agentic Workflows
% ============================================================

\section{Agentic Workflows}

% ---- What is an agent? ------------------------------------
\begin{frame}{What Is an Agent?}
  \begin{columns}[T, onlytextwidth]
    \begin{column}{0.50\textwidth}
      \textbf{Definition}\\[0.3em]
      An \textbf{agent} is an LLM that can:
      \begin{enumerate}
        \item \textbf{Observe} its environment (context, tool outputs)
        \item \textbf{Plan} a sequence of actions
        \item \textbf{Act} by invoking tools
        \item \textbf{Iterate} until the goal is reached
      \end{enumerate}
      \vspace{0.5em}
      This is the \textbf{ReAct} framework~\parencite{yao2023react}: interleaving \emph{reasoning} and \emph{action}.
    \end{column}
    \hfill
    \begin{column}{0.46\textwidth}
      \centering
      \begin{tikzpicture}[
        node distance=0.9cm,
        box/.style={draw, rounded corners=3pt, fill=mLightBrown!15,
                    minimum width=2.0cm, minimum height=0.7cm,
                    align=center, font=\small},
        arr/.style={-{Stealth[length=4pt]}, thick, mLightBrown},
      ]
        \node[box] (obs)  {Observe};
        \node[box, below=of obs]  (plan) {Plan};
        \node[box, below=of plan] (act)  {Act};
        \node[box, below=of act]  (tool) {Tools\\(search, code, file)};
        \node[box, below=of tool] (eval) {Evaluate\\output};

        \draw[arr] (obs.south)  -- (plan.north);
        \draw[arr] (plan.south) -- (act.north);
        \draw[arr] (act.south)  -- (tool.north);
        \draw[arr] (tool.south) -- (eval.north);
        \draw[arr] (eval.east)  -- ++(0.8,0) |- (obs.east)
          node[pos=0.15, right, font=\tiny]{loop};
      \end{tikzpicture}
    \end{column}
  \end{columns}
\end{frame}

% ---- Tool use — building block ----------------------------
\begin{frame}[fragile]{Tool Use — The Building Block}
  \begin{columns}[T, onlytextwidth]
    \begin{column}{0.55\textwidth}
      \textbf{Mechanism}\\[0.3em]
      \begin{itemize}
        \item Model receives a list of available tools as JSON schemas
        \item Model outputs a structured tool-call instead of plain text
        \item Host executes the call, returns results to the model
        \item Model continues reasoning with the result~\parencite{schick2023toolformer}
      \end{itemize}
      \vspace{0.5em}
      \textbf{Common research tools}\\[0.3em]
      \begin{itemize}
        \item \texttt{search\_arxiv(query)} — fetch paper metadata
        \item \texttt{execute\_python(code)} — run analysis code
        \item \texttt{read\_file(path)} / \texttt{write\_file(path, content)}
        \item \texttt{run\_shell(cmd)} — arbitrary commands
        \item \texttt{mcp\_server} — connect to external tools
        \item \textbf{Google Antigravity} — agentic dev platform (IDE)
      \end{itemize}
    \end{column}
    \hfill
    \begin{column}{0.42\textwidth}
      \textbf{Tool schema example}\\[0.3em]
      \begin{lstlisting}[language=json, numbers=none, basicstyle=\ttfamily\scriptsize]
{
  "name": "search_arxiv",
  "description": "Search arXiv for papers",
  "input_schema": {
    "type": "object",
    "properties": {
      "query": {
        "type": "string",
        "description": "Search query"
      },
      "max_results": {
        "type": "integer",
        "default": 5
      }
    },
    "required": ["query"]
  }
}
      \end{lstlisting}
    \end{column}
  \end{columns}
\end{frame}

% ---- Model Context Protocol --------------------------------
\begin{frame}[shrink=18]{Model Context Protocol (MCP)}
  \begin{columns}[T, onlytextwidth]
    \begin{column}{0.55\textwidth}
      \textbf{What is MCP?}~\parencite{anthropic2024mcp}\\[0.3em]
      \begin{itemize}
        \item Open standard for LLM $\leftrightarrow$ external resource connectivity
        \item Proposed by Anthropic; adopted by \tool{Cursor}, \tool{Continue}, \tool{Claude Code}, VS Code, JetBrains
        \item Decouples the \emph{host} (LLM application) from \emph{servers} (data/tool providers)
        \item Communication via JSON-RPC 2.0 over stdio or SSE
        \item No widely-adopted competing open standard; OpenAI's tools API is proprietary
      \end{itemize}
      \vspace{0.5em}
      \textbf{Servers expose three primitives}\\[0.3em]
      \begin{itemize}
        \item \textbf{Resources}: data the model can read (files, DBs)
        \item \textbf{Tools}: functions the model can call
        \item \textbf{Prompts}: reusable prompt templates
      \end{itemize}
    \end{column}
    \hfill
    \begin{column}{0.42\textwidth}
      \centering
      \begin{tikzpicture}[
        node distance=0.8cm and 1.2cm,
        box/.style={draw, rounded corners=2pt, minimum width=2.2cm,
                    minimum height=0.65cm, align=center, font=\tiny},
        arr/.style={-{Stealth[length=4pt]}, thick},
      ]
        \node[box, fill=mLightBrown!20] (host)   {LLM Host\\(Claude Code)};
        \node[box, fill=blue!15, below=of host]   (client) {MCP Client};
        \node[box, fill=green!15, below=of client](server) {MCP Server};
        \node[box, fill=gray!15, below left=0.5cm and -0.4cm of server]  (r1) {Resources};
        \node[box, fill=gray!15, below=0.5cm of server]                  (r2) {Tools};
        \node[box, fill=gray!15, below right=0.5cm and -0.4cm of server] (r3) {Prompts};

        \draw[arr, <->] (host.south)   -- (client.north) node[midway,right,font=\tiny]{};
        \draw[arr, <->] (client.south) -- (server.north) node[midway,right,font=\tiny]{JSON-RPC};
        \draw[arr] (server.south west) -- (r1.north);
        \draw[arr] (server.south)      -- (r2.north);
        \draw[arr] (server.south east) -- (r3.north);
      \end{tikzpicture}
    \end{column}
  \end{columns}
\end{frame}

% ---- Research applications ---------------------------------
\begin{frame}{Research Applications of Agents}
  \begin{columns}[T, onlytextwidth]
    \begin{column}{0.48\textwidth}
      \begin{block}{Literature Sweeps}
        Agent queries arXiv/Semantic Scholar, filters by relevance score, downloads PDFs, extracts structured information, and produces a synthesis report.
      \end{block}
      \vspace{0.4em}
      \begin{block}{Experiment Orchestration}
        Agent reads experiment config, modifies hyperparameters, launches training runs (via \texttt{subprocess}), monitors logs, and emails results.
      \end{block}
    \end{column}
    \hfill
    \begin{column}{0.48\textwidth}
      \begin{block}{Data Pipeline}
        Agent writes a preprocessing script, validates outputs, generates summary statistics, and flags anomalies for human review.
      \end{block}
      \vspace{0.4em}
      \begin{block}{Report Generation}
        Agent reads experiment results, produces tables and figures (matplotlib), and writes a draft results section in LaTeX.
      \end{block}
    \end{column}
  \end{columns}
\end{frame}

% ---- Multi-agent patterns ----------------------------------
\begin{frame}{Multi-Agent Patterns}
  \begin{columns}[T, onlytextwidth]
    \begin{column}{0.50\textwidth}
      \textbf{Orchestrator + Subagents}\\[0.3em]
      \begin{tikzpicture}[
        node distance=0.6cm and 1.2cm,
        box/.style={draw, rounded corners=2pt, minimum width=1.8cm,
                    minimum height=0.6cm, align=center, font=\tiny,
                    fill=mLightBrown!15},
        arr/.style={-{Stealth[length=4pt]}, thick},
      ]
        \node[box] (orch) {Orchestrator LLM};
        \node[box, below left=0.6cm and 0.3cm of orch]  (s1) {Search\\Agent};
        \node[box, below=0.6cm of orch]                  (s2) {Code\\Agent};
        \node[box, below right=0.6cm and 0.3cm of orch] (s3) {Write\\Agent};
        \draw[arr] (orch.south west) -- (s1.north);
        \draw[arr] (orch.south)      -- (s2.north);
        \draw[arr] (orch.south east) -- (s3.north);
        \draw[arr, dashed] (s1.north) -- (orch.south west);
        \draw[arr, dashed] (s2.north) -- (orch.south);
        \draw[arr, dashed] (s3.north) -- (orch.south east);
      \end{tikzpicture}
    \end{column}
    \hfill
    \begin{column}{0.46\textwidth}
      \textbf{Parallel Subagents}\\[0.3em]
      \begin{itemize}
        \item Independent tasks run concurrently
        \item Orchestrator merges results
        \item Reduces total wall-clock time
        \item Example: 10 papers analyzed in parallel
        \item Frameworks: \tool{LangChain}, \tool{AutoGen}, \tool{Google Antigravity}
      \end{itemize}
      \vspace{0.5em}
      \textbf{Reference}\\
      \parencite{wang2024survey}: survey of LLM-based autonomous agents and multi-agent architectures.
    \end{column}
  \end{columns}
\end{frame}

% ---- Harness Engineering Slide ---------------------------
\begin{frame}{Ecosystem Shift: Prompt $\to$ Context $\to$ Harness Engineering}
  \begin{columns}[T, onlytextwidth]
    \begin{column}{0.48\textwidth}
      \textbf{The Harness as Core Object}\\
      The model is only one part of an agent. The \textbf{harness} controls:
      \begin{itemize}
        \item Tool access \& execution permissions
        \item Short-term \& summary memory
        \item Context compaction / token control
        \item Evaluation \& telemetry logging
        \item Multi-agent handoffs / orchestration
      \end{itemize}
      \vspace{0.3em}
      \textbf{Loop Engineering}:
      \begin{itemize}
        \item Run repeated build-test-fix loops autonomously
        \item Automated benchmark execution loops
        \item Specialized peer-reviewer loops
      \end{itemize}
    \end{column}
    \hfill
    \begin{column}{0.48\textwidth}
      \textbf{Foundational Cases / Workflows}
      \begin{itemize}
        \item \textbf{OpenAI's Codex CLI Technical Guidance}: Describes the harness as the agent loop orchestrating user, model, and tools.
        \item \textbf{LangGraph}: Positions itself as an orchestration runtime for streaming, multi-agent coordination, and persistence.
      \end{itemize}
      \vspace{0.4em}
      \begin{block}{Harness over Model}
        \small For research workflows, the harness matters as much as the model: it determines what gets remembered, what tools are safe to call, what gets retried, and what evidence is logged.
      \end{block}
    \end{column}
  \end{columns}
\end{frame}

% ---- Popular agent harnesses ------------------------------
\begin{frame}{Popular Agent Harnesses (2026)}
  \scriptsize
  \begin{tabularx}{\textwidth}{>{\raggedright\arraybackslash}p{0.18\textwidth} >{\raggedright\arraybackslash}p{0.32\textwidth} X}
    \toprule
    \textbf{Harness} & \textbf{Strength} & \textbf{Good fit for research} \\
    \midrule
    \tool{LangGraph} (LC) & Durable graph, persistence, Deep Agents & Complex pipelines with checkpoints, retries, and human approval gates. \\
    \tool{Pi} & Minimal, hackable; packageable skills/extensions & Researchers who want to customize/inspect the loop over using heavy frameworks. \\
    \tool{smolagents} (HF) & Minimal HF library, executable Python code & Quick prototypes and tool loops in a few lines of code. \\
    \tool{OpenAI Agents SDK} & Tool-first hosted agents & Strong when your stack is already centered on OpenAI tools and tracing. \\
    \tool{CrewAI} / AutoGen & Multi-agent conversation/role patterns & Lightweight team-style or role-collaborative (planner, coder, critic) workflows. \\
    \bottomrule
  \end{tabularx}
  \vspace{0.4em}
  \normalsize
  \textbf{Selection heuristic:} pick by \emph{control model} (graph vs minimal loop), \emph{observability}, \emph{human-in-the-loop support}, and \emph{customizability}.
\end{frame}

% ---- Agent Memory Slide ----------------------------------
\begin{frame}{Agent Memory Is Not Just ``Bigger Context''}
  \begin{columns}[T, onlytextwidth]
    \begin{column}{0.48\textwidth}
      \textbf{Three Memory Patterns}\\[0.3em]
      \begin{itemize}
        \item \textbf{Context Compaction}: Dynamically summarizing old interactions or compressing raw text to fit the attention window.
        \item \textbf{Project Memory}: File-based repo context, such as \texttt{CLAUDE.md}, repository maps, local notes, and current task state.
        \item \textbf{Structured Memory Graphs}: e.g., \tool{Beads} (VirtusLab), which provides persistent, structured issue-tracking memory and replaces flat markdown plans with a dependency-aware graph across model restarts.
      \end{itemize}
    \end{column}
    \hfill
    \begin{column}{0.48\textwidth}
      \textbf{Portable Agent Memory}\\[0.3em]
      Citing emerging research proposing structured memory transfer across heterogeneous agents with provenance and fine-grained access control.
      \vspace{0.4em}
      \begin{alertblock}{Warning: The Memory Trade-off}
        Memory increases continuity, but also increases \textbf{privacy risk, stale code assumptions, prompt-injection persistence}, and the \textbf{accidental reuse of wrong decisions}.
      \end{alertblock}
    \end{column}
  \end{columns}
\end{frame}

% ---- Persistent Agents Slide -----------------------------
\begin{frame}{Persistent Agents: Beyond the IDE}
  \small
  \begin{tabularx}{\textwidth}{p{0.22\textwidth} X}
    \toprule
    \textbf{System} & \textbf{Capabilities \& Focus} \\
    \midrule
    \tool{OpenClaw} & Self-hosted personal agent controlled from chat apps; manages email, calendar, files, and tools. Strong data-sovereignty story, but high security risk if over-permissioned (clears inboxes, sends mail). \\
    \midrule
    \tool{Hermes Agent} (Nous Research) & Self-hosted agent with persistent memory and skill creation. An agent that ``grows with you,'' creates new skills from experience, and searches past conversations. \\
    \bottomrule
  \end{tabularx}
  \vspace{0.6em}
  \begin{alertblock}{Key Practice Lesson}
    Persistent operational agents require strict \textbf{permission boundaries, detailed audit logs, memory hygiene}, and \textbf{failsafe revocation controls} to prevent runaway processes or security breaches.
  \end{alertblock}
\end{frame}

% ---- Current limitations -----------------------------------
\begin{frame}{Current Limitations}
  \begin{columns}[T, onlytextwidth]
    \begin{column}{0.48\textwidth}
      \begin{alertblock}{Reliability issues}
        \begin{itemize}
          \item \textbf{Error propagation}: mistakes in early steps compound; agents rarely self-correct without human intervention
          \item \textbf{Cost}: multi-step agents consume many tokens; long tasks can be expensive
        \end{itemize}
      \end{alertblock}
    \end{column}
    \hfill
    \begin{column}{0.48\textwidth}
      \begin{alertblock}{Verification gaps}
        \begin{itemize}
          \item \textbf{Verification}: agents cannot always verify their own outputs; scientific claims must be checked
          \item \textbf{Human-in-the-loop}: for consequential actions (delete files, push to git, email), always require explicit approval
        \end{itemize}
      \end{alertblock}
    \end{column}
  \end{columns}
  \vspace{0.6em}
  \textbf{Best practice:} Use agents for \emph{reversible, low-stakes} tasks. Build in checkpoints for irreversible actions.
  \vspace{0.4em}

  \textbf{Progress note (early 2026):} Coding agents have improved dramatically. SWE-agent v1.07 resolves $\sim$68\% of SWE-bench Verified issues; Claude Opus 4.8 (Thinking) reaches 83.1\%; GPT-5.5 achieves 79.5\% on Terminal-Bench 2.0 and 59.4\% on SWE-Bench Pro. Human oversight remains essential for production code. \textit{(As of June 2026)}
\end{frame}

% ---- Demo idea -------------------------------------------------------
\begin{frame}[shrink=8]{Live Demo Idea: Research Harness on Planning Benchmarks}
  \textbf{Goal (7--8 min):} show one harness coordinating search, coding, and writing in a single loop.

  \vspace{0.4em}
  \textbf{Scenario (from this repo)}
  \begin{itemize}
    \item Start in the Pyperplan workspace (Section 5 context)
    \item Task: compare 2 heuristics on one benchmark domain and produce a mini result summary
  \end{itemize}

  \vspace{0.3em}
  \textbf{Harness roles}
  \begin{itemize}
    \item \textbf{Planner agent}: chooses benchmark/problem pairs and run commands
    \item \textbf{Coder agent}: patches experiment script + parses run logs into a small CSV
    \item \textbf{Writer agent}: drafts 5 bullet points + 1 LaTeX-ready table row
  \end{itemize}

  \vspace{0.3em}
  \textbf{Why it connects to this talk}
  \begin{itemize}
    \item \textbf{Coding assistants}: concrete repo edits and terminal execution
    \item \textbf{Agentic workflows}: orchestrator + subagents + tool calls
    \item \textbf{Writing support}: automatic conversion of results into slide/paper text
  \end{itemize}

  \vspace{0.3em}
  \textbf{Safety gate:} require manual approval before any file overwrite outside a demo folder.
\end{frame}
